%==============================================================================
% Capítulo 7: Consideraciones éticas y aspectos legales (reubicado por c-105)
%==============================================================================
\section{Consideraciones éticas y aspectos legales}\label{sec:aspectos-legales}

\subsection{Marco normativo aplicable}\label{subsec:marco-normativo}

El sistema desarrollado se enmarca normativamente en la Ley~25.326 de Protección de los Datos Personales de la República Argentina \parencite{Ley25326_2000}, en su decreto reglamentario y en las disposiciones complementarias emitidas por la Agencia de Acceso a la Información Pública en su carácter de autoridad de aplicación. Dado que el sistema realiza transferencia de fragmentos textuales hacía infraestructura de inferencia operada por proveedores con sede en los Estados Unidos, resultan aplicables las consideraciones del artículo 12 de la mencionada ley sobre transferencia internacional de datos personales y, de manera complementaria, los lineamientos del Reglamento General de Protección de Datos de la Unión Europea \parencite{GDPR2016} en aquellos casos que pudieran involucrar a residentes europeos. La Unión Europea reconoce a la Argentina como país con nivel adecuado de protección de datos personales mediante la Decisión 2003/490/CE, del 30 de junio de 2003, revalidada en enero de 2024. Esta adecuación ampara el flujo desde la Unión Europea hacia la Argentina y no opera en sentido inverso: los Estados Unidos no revisten la condición de país adecuado para la normativa argentina, de modo que la transferencia desde la Argentina hacia los proveedores de inferencia y de telefonía con sede en ese país requiere un instrumento específico conforme al artículo 12 de la Ley~25.326, según se detalla en la \autoref{subsec:transferencia-internacional}.

\subsection{Principios aplicados al tratamiento de datos}\label{subsec:principios-datos}

El diseño del sistema se atiene de manera explícita a cinco principios fundamentales del tratamiento de datos personales. El principio de licitud establece que todo tratamiento debe ampararse en una base jurídica valida; en este caso, dicha base se encuentra en el artículo 5, inciso 2, apartado d) de la Ley~25.326, que admite el tratamiento de los datos personales que deriven de una relación contractual, científica o profesional del titular y resulten necesarios para su desarrollo o cumplimiento. La organización es la empleadora y trata datos de sus propios empleados para la gestión interna de los incidentes informáticos que estos reportan en ejercicio de sus funciones laborales. Esta base de licitud operativa se distingue del consentimiento que se analiza, para el corpus de investigación académica, en la \autoref{subsec:consentimiento-derechos}, y no dispensa el deber de información del artículo 6 de la Ley~25.326, que subsiste en su integridad. El principio de finalidad determinada limita el uso de los datos a la finalidad declarada de registro y derivación de incidentes, prohibiendo cualquier uso secundario no compatible. El principio de minimización exige recolectar únicamente los datos estrictamente necesarios. En consecuencia, el sistema almacena exclusivamente el nombre de usuario corporativo y un directorio institucional acotado de empleados ---legajo, nombre, correo electrónico corporativo, teléfono, sector y rol---, cuya única finalidad es la resolución de contacto y la notificación al usuario, y descarta cualquier otro dato personal que pudiera capturarse incidentalmente, tales como domicilios particulares, números de documento o información financiera. El principio de exactitud se materializa mediante la posibilidad ofrecida al usuario de revisar y corregir el contenido del incidente antes de su persistencia definitiva. En línea con este principio, el ejercicio de los derechos de rectificación, actualización y supresión (derechos ARCO) se resuelve en un plazo de cinco días hábiles, conforme al artículo 16 de la Ley~25.326. El principio de limitación del plazo de conservación se traduce, para los incidentes, en un plazo de retención de 12 meses rodantes desde la creación del incidente, justificado por el ciclo estadístico anual sobre el que se realiza el análisis de tendencias y la identificación de problemas recurrentes. Este plazo rige para los incidentes y es independiente de la retención de los registros de ejecución del sistema, que se conservan durante 30 días con fines de auditoría y diagnóstico operativo (\autoref{subsec:capa-orquestacion}): se trata de dos conjuntos de datos distintos, con finalidades distintas y plazos propios, de modo que ambas políticas resultan compatibles. Los datos del directorio de empleados, en cambio, se conservan mientras la relación laboral permanezca activa y hasta un año después de su finalización, transcurrido el cual se eliminan, en tanto cumplen una función operativa de contacto y no un fin estadístico. Vencido el plazo de retención, los datos del incidente se someten a anonimización irreversible, de modo que se eliminan los identificadores que permiten vincularlos a una persona. El bloqueo de escritura sobre los incidentes en estado cerrado, que impide editarlos o eliminarlos desde la interfaz de usuario y rechaza con un error 409 los intentos de modificación, constituye un control de integridad de auditoría: preserva la trazabilidad del registro durante el plazo de retención vigente, pero no configura una política de conservación ni una excepción permanente al vencimiento de ese plazo. La pseudonimización de datos personales y el cifrado en reposo garantizan la protección de la información durante el plazo de conservación declarado.

\subsection{Transferencia internacional y pseudonimización}\label{subsec:transferencia-internacional}

El tratamiento de incidentes comporta dos transferencias internacionales de datos hacia proveedores con sede en los Estados Unidos. La primera corresponde al modelo de inferencia Gemini~3.6 Flash, al que se envian fragmentos textuales descriptivos de los incidentes. En este canal, el sistema aplica un procedimiento de pseudonimización dentro del módulo Python desplegado en infraestructura local de la organización, mediante el cual se reemplazan automáticamente los nombres propios, las direcciones de correo electrónico, los números telefónicos, los identificadores corporativos y los nombres de hosts internos por etiquetas genéricas tales como \texttt{[PERSONA]}, \texttt{[EMAIL]} o \texttt{[HOST]}; de este modo, el texto que se transmite no contiene identificadores directos. Esta pseudonimización se valida mediante un conjunto de expresiones regulares específicas del dominio, acompañadas de pruebas unitarias dedicadas.

La segunda transferencia corresponde al canal telefónico, operado por el proveedor Twilio, con sede en los Estados Unidos. El orden real del tratamiento en este canal es: el llamante deja su solicitud y Twilio graba el audio; a continuación, el audio crudo y el número telefónico del llamante se transfieren a la infraestructura del proveedor en los Estados Unidos y se transcriben mediante el servicio de reconocimiento del habla; recién después, el backend aplica la pseudonimización sobre el texto ya transcrito. En consecuencia, el audio crudo del llamante y su número telefónico cruzan a los Estados Unidos antes de cualquier pseudonimización, que opera únicamente sobre el texto y no sobre el audio.

Ambas transferencias se dirigen a un país que no reviste la condición de adecuado para la normativa argentina, por lo que requieren un instrumento específico conforme al artículo 12 de la Ley~25.326: el consentimiento del titular o la suscripción de cláusulas contractuales modelo. El acuerdo de procesamiento de datos suscripto con cada proveedor complementa ese instrumento, pero no lo sustituye, en tanto no configura por si mismo una de las bases de transferencia que la ley admite.

\subsection{Seguridad técnica}\label{subsec:seguridad-tecnica}

Las medidas de seguridad técnica implementadas se sustentan en el modelo de la triada confidencialidad, integridad y disponibilidad (CIA) descrito por \textcite{Stallings2017_computer_security}. La confidencialidad se sustenta, en el diseño, en el cifrado en tránsito sobre TLS~1.2/1.3 a nivel del proxy inverso (Nginx) que protege el perímetro del sistema; el despliegue publicado para la validación experimental no incluyó el proxy, por lo que esa capa de protección corresponde al diseño y no a una capacidad verificada en la instancia publicada. A nivel de aplicación, la confidencialidad se garantiza mediante el cifrado con Fernet (AES-128 en modo CBC con HMAC-SHA-256 integrado, provisto por la biblioteca \texttt{cryptography}) para los campos sensibles del incidente (descripción original, número telefónico del llamante y transcripción). Los datos de contacto del directorio de empleados ---correo corporativo y teléfono--- se almacenan sin cifrado a fin de preservar la auditabilidad directa y las consultas operativas; su protección se sustenta en el control de acceso por rol, el aislamiento de la red interna de Docker y el registro de auditoría de los accesos. La comunicación interna entre los servicios del sistema (N8N, backend, PostgreSQL, Redis) se realiza sobre la red aislada de Docker, que constituye el perímetro de confianza. La integridad se asegura mediante el cifrado autenticado Fernet, que incorpora HMAC-SHA-256 en su modo de operación, la validación estricta de esquemas mediante Pydantic en cada llamada a la interfaz REST, y un registro de auditoría de toda modificación que conserva el actor, la marca temporal y el delta del cambio. La disponibilidad se respalda mediante respaldos diarios automáticos con retención escalonada, despliegue en contenedores con políticas de reinicio automático ante fallas, y monitoreo de salud activo mediante endpoints de chequeo. Las credenciales de servicios externos se gestionan exclusivamente mediante variables de entorno inyectadas por el orquestador de contenedores y nunca se almacenan en el repositorio de código fuente.

\subsection{Consentimiento, derechos del usuario y supervisión humana}\label{subsec:consentimiento-derechos}

Los usuarios de la organización son informados, mediante una política de privacidad accesible desde los canales de entrada, sobre las características del sistema, los datos que se recolectan, la transferencia internacional aplicable y los derechos que les asisten conforme a la normativa vigente. Se garantiza el ejercicio de los derechos de acceso, rectificación, actualización y supresión (derechos ARCO) mediante un procedimiento documentado a cargo del responsable de protección de datos designado por la organización.

En lo que respecta al consentimiento informado para los fines de investigación académica, los 120 empleados de la organización fueron notificados mediante comunicación interna formal, con descripción explícita del propósito del estudio, del tipo de datos tratados, del carácter pseudonimizado de los registros y del derecho a requerir la exclusión de sus incidentes del corpus de validación. El procedimiento fue supervisado por la dirección académica del trabajo y avalado institucionalmente mediante una carta de autorización emitida por la organización adoptante.

En relación con el protocolo de medición en vivo aplicado durante la fase experimental ---en el cual el cronometraje no fue anunciado en tiempo real al operador cuyo tiempo se registraba, con el objeto de mitigar el efecto Hawthorne, pero cuyo conocimiento fue diferido---, corresponde dejar constancia de que, una vez concluida la recolección de datos, se realizó una sesión de esclarecimiento posterior (\textit{debriefing}) con los dos operadores participantes. En dicha sesión se explicó el propósito del estudio, se detalló el procedimiento de medición empleado y se ofreció a cada operador la posibilidad de solicitar la exclusión de sus registros individuales del corpus de validación. Ninguno de los operadores ejerció este derecho. Este procedimiento se alinea con las recomendaciones de la Asociación Americana de Psicología para investigaciones que emplean medición con conocimiento diferido de los participantes \parencite{APA2017}, y fue supervisado por el director del trabajo.

En concordancia con las recomendaciones contemporáneas sobre sistemas que incorporan inteligencia artificial, se preserva en todo momento la posibilidad de supervisión humana significativa. Ninguna decisión que afecte derechos sustanciales de los usuarios se toma de manera completamente automatizada sin posibilidad de revisión, y el sector responsable conserva la potestad de modificar la clasificación inicial cuando lo considere apropiado, registrándose tal modificación en la tabla de trazabilidad descrita en la \autoref{sec:arquitectura}. Este mecanismo materializa el principio de la persona en el bucle (\emph{human in the loop}) como salvaguarda frente a errores potenciales del clasificador automático. Cabe precisar que el umbral de confianza de 0,70 que dispara la revisión obligatoria no constituye, por sí solo, una garantía suficiente de esa salvaguarda: el análisis de calibración de la \autoref{sec:resultados} muestra que la confianza declarada por el clasificador no discrimina de forma fiable los casos correctos de los incorrectos, por lo que la supervisión humana efectiva se apoya además en la revisión por el sector responsable y en la trazabilidad de las correcciones, y no únicamente en el puntaje del modelo.

\newpage
